\subsection*{Question 1}

Consider the portfolio rule:
$
x_{t-1} = \max\left\{0,\; \min\left\{1,\; 
\frac{\mathbb{E}_{t-1}[R_t - R_{f,t}]}{\gamma\sigma^2} \right\}\right\},
$
with risk aversion $\gamma = 5$ and variance $\sigma^2 = 0.0019$.  
The excess portfolio return is
$
r_{p,t} = x_{t-1}(R_t - R_{f,t}),
$
and the annualized Sharpe ratio is $\sqrt{12}$ times the monthly Sharpe ratio.
\subsubsection*{(a) Constant Expected Excess Return}
$$
\mathbb{E}_{t-1}[R_t - R_{f,t}] = 0.0019.
$$
$$
x = \frac{0.0019}{5 \cdot 0.0019} = 0.2.
$$

\textbf{Annualized Sharpe ratio}: $\boxed{\text{SR}_{\text{annual}} = 0.3945}$
\subsubsection*{(b) 5-Year Moving Average Forecast}
$$
\mathbb{E}_{t-1}[R_t - R_{f,t}] = \frac{1}{60}\sum_{s=t-60}^{t-1} (R_s - R_{f,s}),
$$
and $x_{t-1}$ is computed using the same portfolio rule. See the code portion for a list of the expected excess returns.
\\\\
\textbf{Annualized Sharpe ratio}: $\boxed{\text{SR}_{\text{annual}} = 0.2785}$
\subsubsection*{(c) Forecast Based on Dividend Yield}
Estimate the regression using all data prior to $t$:
\[
R_s - R_{f,s} = \alpha_{t-1} + \beta_{t-1} y_{s-1} + \varepsilon_s,
\]
and form the forecast
\[
\mathbb{E}_{t-1}[R_t - R_{f,t}] 
= \alpha_{t-1} + \beta_{t-1} y_{t-1}.
\]
See the code portion for a list of the expected excess returns.
\textbf{Annualized Sharpe ratio}:$\boxed{\text{SR}_{\text{annual}} = 0.3822}$
\subsubsection*{(d) Forecast Based on Net Payout Yield}
Replace $y$ with net payout yield (NPY).
See the code portion for a list of the expected excess returns.
\textbf{Annualized Sharpe ratio}: $\boxed{\text{SR}_{\text{annual}} = 0.3575}$
\subsubsection*{(e) Discussion}
Comparing Sharpe ratios:
\[
\begin{aligned}
\text{Constant mean:} && 0.3945 \\
\text{5-year MA:} && 0.2785 \\
\text{Dividend yield:} && 0.3822 \\
\text{Net payout yield:} && 0.3575
\end{aligned}
\]

The constant‐mean strategy delivers the highest Sharpe ratio (0.3945), while the 5-year moving average performs the worst (0.2785). This occurs because excess returns are extremely noisy, so rolling averages tend to generate unstable and imprecise forecasts.

Dividend yield and net payout yield forecasts do better than the moving average, since these variables are more persistent and theoretically related to expected returns. However, neither predictor produces a Sharpe ratio meaningfully higher than the constant‐mean benchmark. This suggests that although valuation ratios contain some predictive power, the signal is weak relative to return volatility.

Overall, the results show limited evidence of successful mean-based market timing in this sample.
\subsection*{Question 2}
Now consider volatility timing using the rule
$x_{t-1} = 
\max\left\{0,\; 
\min\left\{1,\; 
\frac{E[R_t - R_{f,t}]}{\gamma\sigma_{t-1}^2}
\right\}\right\}
$
with risk aversion $\gamma = 6$ and expected monthly excess return
$
E[R_t - R_{f,t}] = 0.0063.
$ The sample was taken from 1990:2–2014:12.
\subsubsection*{(a) Constant Volatility}
Unconditional variance: $\sigma^2 = 0.001878$

\textbf{Weight}: $x = 0.55896$

\textbf{Annualized Sharpe ratio}: $\boxed{\text{SR}_{\text{annual}} = 0.5224}$

\subsubsection*{(b) Rolling 12-Month Standard Deviation}
See the expected monthly standard deviation in the coding portion.
Consider the equation:
\[
\sigma_{t-1} = \text{Std}(R_{t-12},\ldots,R_{t-1}),
\]
\textbf{Annualized Sharpe ratio}: $\boxed{\text{SR}_{\text{annual}} = 0.5411}$
\subsubsection*{(c) Using VIX}
Monthly volatility estimate:
\[
\sigma_{t-1} = \frac{\text{VIX}_{t-1}}{\sqrt{12}}.
\]

\textbf{Annualized Sharpe ratio}:$\boxed{\text{SR}_{\text{annual}} = 0.6593.}$

\subsubsection*{(d) Discussion}
Sharpe ratio comparison:

\[
\begin{aligned}
\text{Constant volatility:} && 0.5224 \\
\text{Rolling 12-month:} && 0.5411 \\
\text{VIX-based:} && 0.6593
\end{aligned}
\]

Both conditional volatility strategies outperform the constant‐volatility benchmark. The rolling 12-month standard deviation improves performance because volatility is persistent, so scaling down exposure during high‐volatility periods helps stabilize returns.

The VIX-based strategy delivers the highest Sharpe ratio (0.6593). This is intuitive: VIX reflects forward-looking market expectations of volatility, allowing the investor to adjust risk more quickly than with backward-looking realized volatility.

Thus, unlike mean timing in Question 1, volatility timing appears to add meaningful value, especially when using implied volatility.

\pagebreak

\subsection*{Question 3}

\subsubsection*{Setup}
Two investors $i=1,2$, each starts with one unit of wealth and choose holdings $x_i$ of a risky asset (there are two units of risky asset outstanding). Risk-free gross return is $R_f$, risky gross return has mean $\mu$ and variance $\sigma^2$. Investor $i$ receives labor income $Y_i$ with mean $\mu_Y$ and variance $\sigma_Y^2$, and $\operatorname{Cov}(Y_i,R)$ may differ across investors. Define
$$
\beta_i = \frac{\operatorname{Cov}(Y_i,R)}{\sigma^2}.
$$
Each investor maximizes mean--variance utility:
$$
\max_{x_i}\ \ E[W_i] - \frac{\gamma}{2}\operatorname{Var}(W_i),
\qquad
W_i = R_f + x_i(R - R_f) + Y_i.
$$

\subsubsection*{(a) Solve for optimal $x_i$}
Compute $E[W_i]$ and $\operatorname{Var}(W_i)$:
\begin{align*}
E[W_i] &= R_f + x_i(\mu - R_f) + E[Y_i], \\
\operatorname{Var}(W_i) &= x_i^2 \sigma^2 + 2 x_i \operatorname{Cov}(Y_i,R) + \operatorname{Var}(Y_i).
\end{align*}
The FOC w.r.t.\ $x_i$ is
$$
0 = (\mu - R_f) - \gamma\left(x_i \sigma^2 + \operatorname{Cov}(Y_i,R)\right).
$$
Solving:
$$
x_i = \frac{\mu - R_f}{\gamma \sigma^2} - \frac{\operatorname{Cov}(Y_i,R)}{\sigma^2}.
$$
Using $\beta_i$:
$$
\boxed{ \; x_i = \frac{\mu - R_f}{\gamma\sigma^2} - \beta_i \; }.
$$

\subsubsection*{(b) Interpretation}

The optimal position,
\[
x_i = \frac{\mu - R_f}{\gamma\sigma^2} - \beta_i,
\]
has two components. The first term,
\[
x_i^{\text{myopic}} = \frac{\mu - R_f}{\gamma\sigma^2},
\]
is the standard mean--variance demand based solely on expected returns, variance, and risk aversion. The second term,
\[
x_i^{\text{hedge}} = -\beta_i,
\]
adjusts for how investor $i$'s income covaries with the risky asset. If $\beta_i > 0$, income is high when the risky return is high, so the investor reduces risky holdings to hedge this exposure. If $\beta_i < 0$, income is negatively correlated with the risky return, and the investor increases risky holdings to offset income risk. Thus, the hedging term compensates for labor-income risk correlated with the asset.

\subsubsection*{(c) Market clearing and expected excess return}
Market clearing requires
$$
x_1 + x_2 = 2.
$$
Substitute the expression for $x_i$:
$$
\frac{2(\mu - R_f)}{\gamma\sigma^2} - (\beta_1 + \beta_2) = 2.
$$
Solve for $\mu - R_f$:
$$
\boxed{ \; \mu - R_f = \frac{\gamma\sigma^2}{2}\left(2 + \beta_1 + \beta_2\right). \; }
$$

\subsubsection*{(d) Comparative statics on $\beta_2$}

Assume $\beta_1 = 0$. From market clearing, expected excess return satisfies
\[
\mu - R_f = \frac{\gamma\sigma^2}{2}(2 + \beta_2).
\]
If $\beta_2 > 0$, investor 2's income is positively correlated with the risky return, so they reduce their demand for the risky asset. Less demand for risk implies a lower required expected excess return to clear the market.

If $\beta_2 < 0$, investor 2's income is negatively correlated with the risky return, so they demand more of the risky asset to hedge their income risk. Higher demand requires a higher expected excess return to induce sufficient risk-bearing in equilibrium.

Thus, the expected excess return is higher when $\beta_2 < 0$.


\subsection*{Statement of Collaboration} I collaborated with Rosalia Mwidege and Theodore Ouyang. Additionally, ChatGPT was used
to debug any error-prone code and find the proper Excel-equivalent Python
functions and libraries to properly execute the solutions to the problems, in
conjunction with the hints listed on the problem set.

\textbf{Honor Code:} This assignment represents my own work in accordance with
University regulations and class policy.

\end{document}